\section{Prompts for Orchestrator}
\label{app:prompt}



\subsection{Low \masness}

\begin{systemprompt}[title={System Prompt}]
You are a helpful assistant.

\masness (degree of MAS): Low

Valid Channels: thinking, agent, answer

Model: \texttt{[MODEL]}\\

An agent is a pre-configured AI personalities that can delegate tasks to. Each subagent:

1, Has a specific purpose and expertise area\\
2. Uses its own context window separate from the main conversation\\
3. (Optional) Can be configured with specific tools it's allowed to use\\
4. Includes a custom system prompt that guides its behavior \\
\\
An agent can only call within the \xmltag{thinking} channel, it should be defined in tag \xmltag{agent}. Each agent must contain \xmltag{agent_name}, \xmltag{agent_description}, \xmltag{required_arguments}, \xmltag{agent_output_id}\\

DO NOT MISS ANY REQUEST FIELDS and ensure that your response is a well-formed XML object!
\end{systemprompt}


\begin{developprompt}[title={Develop Prompt}]
Channels: \\
    \xmltag{thinking}: internal reasoning and planning \\
    \xmltag{agent}: definition of sub-agents \\
    \xmltag{answer}: final user-facing answer \\

Model (the model used in sub-agent):

    \texttt{gpt-4.1-nano}: [Introduction of the model from official website] \\
    \texttt{gpt-oss-120b}:[Introduction of the model from official website] \\
    (...list of all candidate models)\\

\masness Levels:

    \texttt{low}: direct solve or at most one agent \\
    \texttt{high}: complex multi-agent delegation\\

Sub-agent Schema (all fields required):\\

\xmltag{agent}

\quad\xmltag{agent_name}
...
\xmlendtag{agent_name}
(select one of the agents: CoTAgent, SCAgent, DebateAgent, ReflexionAgent)

\quad\xmltag{agent_description}
...
\xmlendtag{agent_description}

\quad\xmltag{required_arguments}
(make sure all required parameters are set. Must follow XML format)

\qquad\xmltag{...}...\xmlendtag{...}

\qquad\xmltag{...}...\xmlendtag{...}

\quad\xmlendtag{required_arguments}

\quad\xmltag{agent_output_id}
...
\xmlendtag{agent_output_id}

\xmlendtag{agent} \\

Available Agents: \\

\texttt{CoT}: [Description, Name, Required Argument and Summary of Implementation]

\texttt{SC}: [Description, Name, Required Argument and Summary of Implementation]

\texttt{Debate}: [Description, Name, Required Argument and Summary of Implementation]

\texttt{Reflexion}: [Description, Name, Required Argument and Summary of Implementation]

\texttt{Search}: [Description, Name, Required Argument and Summary of Implementation]


\end{developprompt}

\begin{userprompt}[title={User Prompt}]
Please solve the question step-by-step. During the reasoning process, you can address the task yourself and output the final answer in the \xmltag{answer} tag. You can also delegate the task to the agent that you designed, and output the corresponding \texttt{agent_output_id} in the \xmltag{answer} tag. You must output ALL required parameters and use EXACTLY the same field names for the agent. \\
\\
For example,
\\
If you can solve the task yourself, you will output the following:\\

\xmltag{thinking}

(20+9)*(30+7) = 600 + 140 + 270 + 63 = 1073.

\xmlendtag{thinking}

\xmltag{answer}1073\xmlendtag{answer}
\\

If you can solve the task via single-agent delegation with one tool call, you will output the following:\\

\xmltag{thinking}

This problem involves symbolic integration and applying the Fundamental Theorem of Calculus.
It requires structured reasoning rather than simple numeric computation.
I will use a calculus agent that can perform step-by-step Chain-of-Thought reasoning.
The final answer to the original question will be the output of the CoTAgent.

\xmlendtag{thinking}

\xmltag{agent}

\quad\xmltag{agent_name}CoTAgent\xmlendtag{agent_name}

\quad\xmltag{agent_description}{Definite integrals with one Chain-of-Thought call.}\xmlendtag{agent_description}

\quad\xmltag{required_arguments}

\qquad\xmltag{agent_input}\xmlendtag{agent_input}

\quad\xmlendtag{required_arguments}

\quad\xmltag{agent_output_id}
calc_agent_output
\xmlendtag{agent_output_id}

\xmlendtag{agent}

\xmltag{answer}{calc_agent_output}\xmlendtag{answer} \\

Another example:\\

\xmltag{thinking}

This question requires comparison between two close numeric choices.
To ensure accuracy, I will let two reasoning roles debate: one focusing on mathematical precision and the other on practical rounding.
The DebateAgent can capture both perspectives and reach a justified final answer.
The final answer will be the output of the DebateAgent.

\xmlendtag{thinking}

\xmltag{agent}

\quad\xmltag{agent_name}DebateAgent\xmlendtag{agent_name}

\quad\xmltag{agent_description}
Near-tie numeric choice using one Debate call.
\xmlendtag{agent_description}

\quad\xmltag{required_arguments}

\qquad\xmltag{agent_input}\xmlendtag{agent_input}

\qquad\xmltag{debate_roles}
["Mathematics Professor", "Statistics Teacher"]
\xmlendtag{debate_roles}

\quad\xmlendtag{required_arguments}

\quad\xmltag{agent_output_id}
compare_agent_output
\xmlendtag{agent_output_id}

\xmlendtag{agent}

\xmltag{answer}compare_agent_output\xmlendtag{answer}\\

More examples:\\

\xmltag{thinking}

Computing 17³ can be done directly, but arithmetic mistakes are easy to make.
Using a Chain-of-Thought with Self-Consistency (CoT\_SC) allows sampling multiple reasoning paths and combining results for accuracy.
The final answer to the original question will be the output of the SCAgent.

\xmlendtag{thinking}

\xmltag{agent}

\quad\xmltag{agent_name}SCAgent\xmlendtag{agent_name}

\quad\xmltag{agent_description}
Performs arithmetic calculations using Chain-of-Thought with Self-Consistency (CoT\_SC).
\xmlendtag{agent_description}

\quad\xmltag{required_arguments}

\qquad\xmltag{agent_input}\xmlendtag{agent_input}

\quad\xmlendtag{required_arguments}

\quad\xmltag{agent_output_id}math_agent_output\xmlendtag{agent_output_id}

\xmlendtag{agent}

\xmltag{answer}math_agent_output\xmlendtag{answer}
\\

Final example: \\

\xmltag{thinking}

This task requires reasoning with a formula and ensuring units are handled correctly.
I will use a Reflexion agent that can reflect on and refine its reasoning if errors occur.
The final answer to the original question will be the output of the ReflexionAgent.

\xmlendtag{thinking}

\xmltag{agent}

\quad\xmltag{agent_name}ReflexionAgent\xmlendtag{agent_name}

\quad\xmltag{agent_description}
Solves reasoning tasks with iterative self-reflection and critique using Reflexion.
\xmlendtag{agent_description}

\quad\xmltag{required_arguments}

\qquad\xmltag{agent_input}\xmlendtag{agent_input}

\quad\xmlendtag{required_arguments}

\quad\xmltag{agent_output_id}reflexion_agent_output\xmlendtag{agent_output_id}

\xmlendtag{agent}

\xmltag{answer}reflexion_agent_output\xmlendtag{answer} \\


If you combine one direct step with a single Reflexion tool call, you will output the following:

\xmltag{thinking}

Step 1: Compute the average manually:
\quad sum = 2 + 3 + 5 + 7 + 11 = 28
\quad avg = 28 / 5 = 5.6

Step 2:
\quad The remaining step — computing sqrt(5.6) to three decimals — requires precision and numeric refinement.
\quad I will delegate that part to a ReflexionAgent using the ReflexionAgent for self-correction if rounding is wrong.
\quad The final answer to the original question will be the output of the ReflexionAgent.

\xmlendtag{thinking}

\xmltag{agent}

\quad\xmltag{agent_name}ReflexionAgent\xmlendtag{agent_name}

\quad\xmltag{agent_description}
Square root with a light self-refine loop (single agent call).
\xmlendtag{agent_description}

\quad\xmltag{required_arguments}

\qquad\xmltag{agent_input}Compute sqrt(5.6) to 3 decimals\xmlendtag{agent_input}

\quad\xmlendtag{required_arguments}

\quad\xmltag{agent_output_id}numeric_agent_output\xmlendtag{agent_output_id}

\xmlendtag{agent}

\xmltag{answer}numeric_agent_output\xmlendtag{answer} \\


Below is the question to solve: 

[QUESTION]


\end{userprompt}


\subsection{High \masness}

% \zixuan{TODO: need to mention the difference}


\begin{systemprompt}[title={System Prompt}]
You are a helpful assistant.

\masness (degree of MAS): High

Valid Channels: thinking, agent, edge

Model: \texttt{[MODEL]}\\

An agent is a pre-configured AI personalities that can delegate tasks to. Each subagent:
1. Has a specific purpose and expertise area\\
2. Uses its own context window separate from the main conversation\\
3. (Optional) Can be configured with specific tools it's allowed to use\\
4. Includes a custom system prompt that guides its behavior\\
\\
An agent can only call within the \xmltag{thinking} channel, it should be defined in tag \xmltag{agent}. Each agent must contain \xmltag{agent_id}, \xmltag{agent_name}, \xmltag{agent_description}, \xmltag{required_arguments}. To connect multiple agents to form a multi-agent system, use \xmltag{edge} channel.\\


DO NOT MISS ANY REQUEST FIELDS and ensure that your response is a well-formed XML object!
\end{systemprompt}


\begin{developprompt}[title={Develop Prompt}]
Channels: \\
    \xmltag{thinking}: internal reasoning and planning \\
    \xmltag{agent}: definition of sub-agents \\
    \xmltag{answer}: final user-facing answer \\

Model (the model used in sub-agent):

    \texttt{gpt-4.1-nano}: [Introduction of the model from official website] \\
    \texttt{gpt-oss-120b}:[Introduction of the model from official website] \\
    (...list of all candidate models)\\

\masness Levels:

    \texttt{low}: direct solve or at most one agent \\
    \texttt{high}: complex multi-agent delegation\\
    
Sub-agent Schema (all fields required):\\

\xmltag{agent}

\quad\xmltag{agent_name}
...
\xmlendtag{agent_name}
(select one of the agents: CoTAgent, SCAgent, DebateAgent, ReflexionAgent)

\quad\xmltag{agent_description}
...
\xmlendtag{agent_description}

\quad\xmltag{required_arguments}
(make sure all required parameters are set. Must follow XML format)

\qquad\xmltag{...}...\xmlendtag{...}

\qquad\xmltag{...}...\xmlendtag{...}

\quad\xmlendtag{required_arguments}

\quad\xmltag{agent_output_id}
...
\xmlendtag{agent_output_id}

\xmlendtag{agent} \\


Edge Schema (single block; all fields required. Each pair defines a directed link: output of \xmltag{from} → input of \xmltag{to}. List ALL links here; use exactly one \xmltag{edge} block per solution:\\

\xmltag{edge}

\quad\xmltag{from}
...
\xmlendtag{from}
(the source agent\_id)

\quad\xmltag{to}
...
\xmlendtag{to}
(the target agent\_id)

\xmlendtag{edge}\\

Available Agents: \\

\texttt{CoT}: [Description, Name, Required Argument and Summary of Implementation]

\texttt{SC}: [Description, Name, Required Argument and Summary of Implementation]

\texttt{Debate}: [Description, Name, Required Argument and Summary of Implementation]

\texttt{Reflexion}: [Description, Name, Required Argument and Summary of Implementation]

\texttt{Search}: [Description, Name, Required Argument and Summary of Implementation]


\end{developprompt}



\begin{userprompt}[title={User Prompt}]

Please solve the given question by creating one or more agents and connecting them into a valid computational graph that collaboratively produces the final answer. To create an agent, you must define that agent by outputting \xmltag{agent} with \xmltag{agent_id} (a unique id for the agent, must be unique and contain only alphanumeric or underscore characters (e.g., A1, Refine_1, WS_Japan)), \xmltag{agent_name} (exactly one of: CoTAgent, SCAgent, DebateAgent, ReflexionAgent and WebSearchAgent), \xmltag{agent_description}, \xmltag{required_arguments} (must include at least one \xmltag{agent_input} tag. DebateAgents must define \xmltag{debate_roles} with two or more roles. If  \xmltag{agent_input} left empty (""), the parser will automatically replace it with the original question.). \\

After defining all agents, you must build a valid graph by specifying edges that describe the data flow between agents. Output exactly one \xmltag{edge} block, Each \xmltag{from} \xmltag{to} pair connects the output of one agent to the input of another: \\

\xmltag{edge}

\quad\xmltag{from}source_agent_id\xmlendtag{from}

\quad\xmltag{to}target_agent_id\xmlendtag{to}

\xmlendtag{edge} \\

You can output multiple \xmltag{from} and \texttt{\xmltag{/to}} inside the \xmltag{edge}. Each \xmltag{from} and \xmltag{to} value must exactly match an existing \xmltag{agent_id}. \\

To be valid, your graph must satisfy all of the following constraints: 
\begin{itemize}[leftmargin=*,noitemsep,topsep=2pt]
\item Node consistency: Every \xmltag{from} and \xmltag{to} must reference a valid \xmltag{agent_id} that appears in an \xmltag{agent} block.
\item  Directionality: Edges are directed: data flows from \xmltag{from} → \xmltag{to}.
\item  Connectivity: Every agent must be connected directly or indirectly to the main flow. Isolated agents are not allowed.
\item  Start node(s): At least one agent must have no incoming edge. These are “entry points” (e.g., WebSearch or initial reasoning).
\item  Sink node: There must be exactly one agent with no outgoing edge — this is the FINAL agent that produces the answer.
\item  No undefined edges: It is invalid to reference an agent in \xmltag{from} or \xmltag{to} that was not declared.
\item  No loops or cycles: No self-loop: \xmltag{from}X\xmltag{/from}\xmltag{to}X\xmltag{/to} is not allowed. No cycles: The graph must be acyclic; a topological order must exist.
\item  Parallelism allowed: Multiple agents may have the same \xmltag{from} or \xmltag{to} (fan-out/fan-in).
\item  Unambiguous sink: The parser will reject graphs with multiple sinks (add a final “collector” agent if needed).
\item  Order-independent: The XML order of edges does not need to follow execution order; topological sorting is handled automatically.
\item  Sink answer completeness: The unique sink agent's output must directly answer the original question in a user-ready form. It must not be an intermediate artifact (e.g., notes, critique, raw table) unless the question explicitly asks for that artifact. If \xmltag{agent_input} is empty for the sink, it inherits the original question and must return the final answer. If \xmltag{agent_input} is non-empty, the runner still prepends the original question as context; the sink must still produce the final, user-facing answer to that question.
\item Edge-Data Flow Consistency (BIDIRECTIONAL): Edges represent execution order. \$\{\} represents data flow. As a result, you must ensure they are consistent with each other.
    (a) If an \xmltag{agent_input} references \$\{X\}, there MUST be an edge \xmltag{from}X\xmltag{/from}\xmltag{to}THIS AGENT\xmltag{/to}
    (b) If there is an edge \xmltag{from}X\xmltag{/from}\xmltag{to}Y\xmltag{/to}, then Y's \xmltag{agent_input} MUST reference \$\{X\}
    In other words: edges exist if and only if there is actual data passing from one agent to another. Do not create edges solely for execution ordering without data flow.
\end{itemize}

Thinking Section (Required):\\

Before defining agents and edges, you must include a \xmltag{thinking} section.
It should naturally describe why multiple agents are needed, why each type was chosen, and why the graph has that structure (parallel, sequential or hybrid).
It must justify both planning and design rationale.\\

Example structure:\\

\xmltag{thinking}
  Explain why a single agent is insufficient.
  Describe each agent's role and how they connect.
  Justify the flow pattern (parallel, sequential, hybrid).
  End by clearly stating which agent produces the final output.
\xmltag{/thinking}\\

Single-agent example:\\

If you decide to solve via single agent, you will output the following. In this case, since the \xmltag{agent_input} is the same as the original task, you must set the \xmltag{agent_input} as empty (""), and the parser will replace it with the original question.\\

1st example:\\

Question: Compute the definite integral of (2x + 5) dx from 0 to 3.\\

\xmltag{thinking}

This problem involves symbolic integration and applying the Fundamental Theorem of Calculus.
It requires structured reasoning rather than simple numeric computation.
I will use a calculus agent that can perform step-by-step Chain-of-Thought reasoning.
The final answer to the original question will be the output of the CoTAgent.

\xmlendtag{thinking}

\xmltag{agent}

\quad\xmltag{agent_id}calc\_agent\xmlendtag{agent_id}

\quad\xmltag{agent_name}CoTAgent\xmlendtag{agent_name}

\quad\xmltag{agent_description}
Definite integrals with one Chain-of-Thought call.
\xmlendtag{agent_description}

\quad\xmltag{required_arguments}

\qquad\xmltag{agent_input}\xmlendtag{agent_input}

\quad\xmlendtag{required_arguments}

\xmlendtag{agent}\\

2nd example:\\

Question: What is 17 cubed?\\

\xmltag{thinking}

Computing 17³ can be done directly, but arithmetic mistakes are easy to make.
Using a Chain-of-Thought with Self-Consistency (CoT\_SC) allows sampling multiple reasoning paths and combining results for accuracy.
The final answer to the original question will be the output of the SCAgent.

\xmlendtag{thinking}

\xmltag{agent}

\quad\xmltag{agent_id}math\_agent\xmlendtag{agent_id}

\quad\xmltag{agent_name}SCAgent\xmlendtag{agent_name}

\quad\xmltag{agent_description}
Performs arithmetic calculations using Chain-of-Thought with Self-Consistency (CoT\_SC).
\xmlendtag{agent_description}

\quad\xmltag{required_arguments}

\qquad\xmltag{agent_input}\xmlendtag{agent_input}

\quad\xmlendtag{required_arguments}

\xmlendtag{agent} \\

3rd examples: \\

Question: Given x² = 46.694444, which target number is closer, 45 or 46? \\

\xmltag{thinking}

This question requires comparison between two close numeric choices.
To ensure accuracy, I will let two reasoning roles debate — one focusing on mathematical precision and the other on practical rounding.
The DebateAgent can capture both perspectives and reach a justified final answer.
The final answer will be the output of the DebateAgent.

\xmlendtag{thinking}

\xmltag{agent}

\quad\xmltag{agent_id}compare\_agent\xmlendtag{agent_id}

\quad\xmltag{agent_name}DebateAgent\xmlendtag{agent_name}

\quad\xmltag{agent_description}
Near-tie numeric choice using one Debate call.
\xmlendtag{agent_description}

\quad\xmltag{required_arguments}

\qquad\xmltag{agent_input}\xmlendtag{agent_input}

\qquad\xmltag{debate_roles}
["Mathematics Professor", "Statistics Teacher"]
\xmlendtag{debate_roles}

\quad\xmlendtag{required_arguments}

\xmlendtag{agent}

4th example: \\

Question: A train travels 60 miles per hour. How far does it go in 2.5 hours? \\

\xmltag{thinking}

This task requires reasoning with a formula and ensuring units are handled correctly.
I will use a Reflexion agent that can reflect on and refine its reasoning if errors occur.
The final answer to the original question will be the output of the ReflexionAgent.

\xmlendtag{thinking}

\xmltag{agent}

\quad\xmltag{agent_id}reflexion\_agent\xmlendtag{agent_id}

\quad\xmltag{agent_name}ReflexionAgent\xmlendtag{agent_name}

\quad\xmltag{agent_description}
Solves reasoning tasks with iterative self-reflection and critique using Reflexion.
\xmlendtag{agent_description}

\quad\xmltag{required_arguments}

\qquad\xmltag{agent_input}\xmlendtag{agent_input}

\quad\xmlendtag{required_arguments}

\xmlendtag{agent}\\

5th example:\\

Question: What is the current inflation rate in Japan as of this month?\\

\xmltag{thinking}

This question depends on up-to-date factual information that cannot be reliably recalled from static knowledge.  
Using a single WebSearchAgent is sufficient because the task only requires retrieving accurate, cited facts from the web, not further reasoning or synthesis.  
The agent will search online sources and return a concise, citation-based summary.  
The WebSearchAgent is therefore both the only node and the final output of the flow.

\xmlendtag{thinking}

\xmltag{agent}

\quad\xmltag{agent_id}SEARCH\xmlendtag{agent_id}

\quad\xmltag{agent_name}WebSearchAgent\xmlendtag{agent_name}

\quad\xmltag{agent_description}
Retrieves recent and cited factual information from the internet.
\xmlendtag{agent_description}

\quad\xmltag{required_arguments}

\qquad\xmltag{agent_input}\xmlendtag{agent_input}

\quad\xmlendtag{required_arguments}

\xmlendtag{agent}\\

Multi-Agent Example: \\

If you decide to solve via multiple agents, you must first decompose the original question into smaller, well-defined sub-tasks, each representing a single sub-goal. Then, create one agent per sub-task and connect them into a coherent, acyclic computational graph.\\

In this case, the agent_input is not empty and serves as the specific sub-task for the agent to solve, while the parser automatically prepends the original question as context before the provided agent_input content. \\

When decomposing:\\

1. Keep sub-tasks minimal and focused. Each agent should handle one atomic objective (e.g., one query, one reasoning step, or one verification task). \\
2. Use multiple WebSearchAgents for different pieces of factual evidence, rather than a single broad search.\\
3. Connect agents logically so that information flows toward a single final agent (the sink) that directly answers the original question.\\


1st example:\\

Question: Give a short, cited summary of the most recent housing vacancy rates for New York City, Los Angeles, and Chicago\\


\xmltag{thinking}

We need up-to-date numbers with sources, so we will run three WebSearch agents in parallel.
Then a CoT agent will normalize the three snippets into a small table.
An SC agent will run small variants and vote to reduce extraction errors.
A Reflexion agent will check units, recency, and citations.
A final CoT agent will write the short summary.
The final sink is FINAL.

\xmlendtag{thinking}

\xmltag{agent}

\quad\xmltag{agent_id}WS\_NYC\xmlendtag{agent_id}

\quad\xmltag{agent_name}WebSearchAgent\xmlendtag{agent_name}

\quad\xmltag{agent_description}
Find the latest NYC housing vacancy rate with source text.
\xmlendtag{agent_description}

\quad\xmltag{required_arguments}

\qquad\xmltag{agent_input}
Search for the latest official or reputable estimate of the housing vacancy rate for New York City. Return a short snippet with the number, date, and a citation line.
\xmlendtag{agent_input}

\quad\xmlendtag{required_arguments}

\xmlendtag{agent}

\xmltag{agent}

\quad\xmltag{agent_id}WS\_LA\xmlendtag{agent_id}

\quad\xmltag{agent_name}WebSearchAgent\xmlendtag{agent_name}

\quad\xmltag{agent_description}
Find the latest LA vacancy rate with source text.
\xmlendtag{agent_description}

\quad\xmltag{required_arguments}

\qquad\xmltag{agent_input}
Search for the latest official or reputable estimate of the housing vacancy rate for Los Angeles. Return a short snippet with the number, date, and a citation line.
\xmlendtag{agent_input}

\quad\xmlendtag{required_arguments}

\xmlendtag{agent}

\xmltag{agent}

\quad\xmltag{agent_id}WS\_CHI\xmlendtag{agent_id}

\quad\xmltag{agent_name}WebSearchAgent\xmltag{agent_name}

\quad\xmltag{agent_description}
Find the latest Chicago vacancy rate with source text.
\xmlendtag{agent_description}

\quad\xmltag{required_arguments}

\qquad\xmltag{agent_input}
Search for the latest official or reputable estimate of the housing vacancy rate for Chicago. Return a short snippet with the number, date, and a citation line.
\xmlendtag{agent_input}

\quad\xmlendtag{required_arguments}

\xmlendtag{agent}

\xmltag{agent}

\quad\xmltag{agent_id}EXT\xmlendtag{agent_id}

\quad\xmltag{agent_name}CoTAgent\xmlendtag{agent_name}

\quad\xmltag{agent_description}
Extract numbers and standardize the three rates with dates and citations.
\xmlendtag{agent_description}

\quad\xmltag{required_arguments}

\qquad\xmltag{agent_input}
From the snippets below, extract for each city: city name, vacancy rate (as a percent), reference date (YYYY-MM or YYYY), and a short citation. Output a compact 3-line table.

NYC:
\$\{WS\_NYC\}

LA:
\$\{WS\_LA\}

Chicago:
\$\{WS\_CHI\}
\xmlendtag{agent_input}

\quad\xmlendtag{required_arguments}

\xmlendtag{agent}

\xmltag{agent}

\quad\xmltag{agent_id}VOTE\xmlendtag{agent_id}

\quad\xmltag{agent_name}SCAgent\xmlendtag{agent_name}

\quad\xmltag{agent_description}
Ensemble the extraction to reduce copy or parse errors.
\xmlendtag{agent_description}

\quad\xmltag{required_arguments}

\qquad\xmltag{agent_input}
Given the 3-line table below, produce 5 independent extractions and vote on a single corrected 3-line table.

Table:
\$\{EXT\}
\xmlendtag{agent_input}

\quad\xmlendtag{required_arguments}

\xmlendtag{agent}

\xmltag{agent}

\quad\xmltag{agent_id}QA\xmlendtag{agent_id}

\quad\xmltag{agent_name}ReflexionAgent\xmlendtag{agent_name}

\quad\xmltag{agent_description}
Check units, date freshness, and cite presence; list fixes if needed.
\xmlendtag{agent_description}

\quad\xmltag{required_arguments}

\qquad\xmltag{agent_input}
Audit the voted table for: units as \%, dates present, and a citation per city. If any issue is found, list concrete fixes in 3 lines; else say OK. End with a 1-line verdict.

Voted table:
\$\{VOTE\}
\xmlendtag{agent_input}

\quad\xmlendtag{required_arguments}

\xmlendtag{agent}

\xmltag{agent}

\quad\xmltag{agent_id}FINAL\xmlendtag{agent_id}

\quad\xmltag{agent_name}CoTAgent\xmlendtag{agent_name}

\quad\xmltag{agent_description}
Write the short, cited summary.
\xmlendtag{agent_description}

\quad\xmltag{required_arguments}

\qquad\xmltag{agent_input}
Using the checked table and notes, write a 3--4 sentence summary with one sentence per city and a final sentence comparing the rates. Keep citations as inline source lines from the table.

Table:
\$\{VOTE\}

QA notes:
\$\{QA\}
\xmlendtag{agent_input}

\quad\xmlendtag{required_arguments}

\xmlendtag{agent}

\xmltag{edge}

\quad\xmltag{from}WS\_NYC\xmlendtag{from}\xmltag{to}EXT\xmlendtag{to}

\quad\xmltag{from}WS\_LA\xmlendtag{from}\xmltag{to}EXT\xmlendtag{to}

\quad\xmltag{from}WS\_CHI\xmlendtag{from}\xmlendtag{from}\xmltag{to}EXT\xmlendtag{to}

\quad\xmltag{from}EXT\xmlendtag{from}\xmltag{to}VOTE\xmlendtag{to}

\quad\xmltag{from}VOTE\xmlendtag{from}\xmltag{to}QA\xmlendtag{to}

\quad\xmltag{from}VOTE\xmlendtag{from}\xmltag{to}FINAL\xmlendtag{to}

\quad\xmltag{from}QA\xmlendtag{from}\xmltag{to}FINAL\xmlendtag{to}

\xmlendtag{edge}\\


2nd example:\\

Question: During Pope John Paul II's first foreign journey in the late 1970s, he visited a country known for its rich Mesoamerican history and home to a large population. On which other date did he visit a major city on the Adriatic Sea, known for its significant port and a famous basilica dedicated to a saint with the initial “S”, and which other nearby city did he visit on the same day?\\


\xmltag{thinking}

The question needs historical reasoning: identify Pope John Paul II's first foreign trip, find an Adriatic city he visited with a basilica of a saint starting with “S”, get the date, and determine another city visited the same day.

One agent cannot do all of this because it mixes retrieval and reasoning.
I will decompose it into smaller sub-tasks: three WebSearchAgents for each factual lookup, one CoTAgent to build a timeline, one ReflexionAgent to verify consistency, and a final CoTAgent to write the answer.
The FINAL agent outputs the final answer.

\xmlendtag{thinking}

\xmltag{agent}

\quad\xmltag{agent_id}WS\_FIRST\_TRIP\xmlendtag{agent_id}

\quad\xmltag{agent_name}WebSearchAgent\xmlendtag{agent_name}

\quad\xmltag{agent_description}
Retrieve Pope John Paul II's first foreign trip details in the late 1970s.
\xmlendtag{agent_description}

\quad\xmltag{required_arguments}

\qquad\xmltag{agent_input}
Search for Pope John Paul II's first foreign journey (year, destination country, and duration).
Return the trip date range, destination, and a reliable citation.
\xmlendtag{agent_input}

\quad\xmlendtag{required_arguments}

\xmlendtag{agent}

\xmltag{agent}

\quad\xmltag{agent_id}WS\_ADRIATIC\xmlendtag{agent_id}

\quad\xmltag{agent_name}WebSearchAgent\xmlendtag{agent_name}

\quad\xmltag{agent_description}
Find Adriatic Sea city visits with basilicas dedicated to saints starting with `S'.
\xmlendtag{agent_description}

\quad\xmltag{required_arguments}

\qquad\xmltag{agent_input}
Search for any Adriatic city visited by Pope John Paul II that has a basilica dedicated to a saint with the initial `S'
(e.g., Saint Nicholas, Saint Mark).
Return the city name, basilica name, and visit date, with citation.
\xmlendtag{agent_input}

\quad\xmlendtag{required_arguments}

\xmlendtag{agent}

\xmltag{agent}

\quad\xmltag{agent_id}WS\_SAME\_DAY\xmlendtag{agent_id}

\quad\xmltag{agent_name}WebSearchAgent\xmlendtag{agent_name}

\quad\xmltag{agent_description}
Identify any other nearby city visited by Pope John Paul II on the same day as the Adriatic visit.
\xmlendtag{agent_description}

\quad\xmltag{required_arguments}

\qquad\xmltag{agent_input}
Search for other cities visited by Pope John Paul II on the same date as his Adriatic visit.
Return the nearby city name, distance estimate, and source citation.
\xmlendtag{agent_input}

\quad\xmlendtag{required_arguments}

\xmlendtag{agent}

\xmltag{agent}

\quad\xmltag{agent_id}TIMELINE\xmlendtag{agent_id}

\quad\xmltag{agent_name}CoTAgent\xmlendtag{agent_name}

\quad\xmltag{agent_description}
Integrate trip data into a single verified historical timeline.
\xmlendtag{agent_description}

\quad\xmltag{required_arguments}

\qquad\xmltag{agent_input}
Using the results below:

First foreign journey:
\$\{WS\_FIRST\_TRIP\}

Adriatic visit:
\$\{WS\_ADRIATIC\}

Same-day visit:
\$\{WS\_SAME\_DAY\}

Build a clear timeline confirming the Adriatic visit date and same-day nearby city.
\xmlendtag{agent_input}

\quad\xmlendtag{required_arguments}

\xmlendtag{agent}

\xmltag{agent}

\quad\xmltag{agent_id}VERIFY\xmlendtag{agent_id}

\quad\xmltag{agent_name}ReflexionAgent\xmlendtag{agent_name}

\quad\xmltag{agent_description}
Validate chronology, geography, and saint-basilica link.
\xmlendtag{agent_description}

\quad\xmltag{required_arguments}

\qquad\xmltag{agent_input}
Check consistency among timeline facts:
- Ensure the same date appears across all sources.
- Confirm the Adriatic city is geographically near the second city.
- Verify that the basilica indeed honors a saint whose name starts with `S'.

Return OK if consistent, otherwise list corrections.

Timeline to verify:
\$\{TIMELINE\}
\xmlendtag{agent_input}

\quad\xmlendtag{required_arguments}

\xmlendtag{agent}

\xmltag{agent}

\quad\xmltag{agent_id}FINAL\xmlendtag{agent_id}

\quad\xmltag{agent_name}CoTAgent\xmlendtag{agent_name}

\quad\xmltag{agent_description}
Produce the final concise answer with the date and both cities.
\xmlendtag{agent_description}

\quad\xmltag{required_arguments}

\qquad\xmltag{agent_input}
Using verified results, answer in one sentence:

``Pope John Paul II visited [ADRIATIC\_CITY] on [DATE], home to the Basilica of Saint [S], and also visited [NEARBY\_CITY] on the same day.''

Include a brief verification sentence citing sources.

Timeline:
\$\{TIMELINE\}

Verification:
\$\{VERIFY\}
\xmlendtag{agent_input}

\quad\xmlendtag{required_arguments}

\xmlendtag{agent}

\xmltag{edge}

\quad\xmltag{from}WS\_FIRST\_TRIP\xmlendtag{from}\xmltag{to}TIMELINE\xmlendtag{to}

\quad\xmltag{from}WS\_ADRIATIC\xmlendtag{from}\xmltag{to}TIMELINE\xmlendtag{to}

\quad\xmltag{from}WS\_SAME\_DAY\xmlendtag{from}\xmltag{to}TIMELINE\xmlendtag{to}

\quad\xmltag{from}TIMELINE\xmlendtag{from}\xmltag{to}VERIFY\xmlendtag{to}

\quad\xmltag{from}VERIFY\xmlendtag{from}\xmltag{to}FINAL\xmlendtag{to}

\quad\xmltag{from}TIMELINE\xmlendtag{from}\xmltag{to}FINAL\xmlendtag{to}

\xmlendtag{edge}\\

(... more examples) \\

Below is the question to solve: 

[QUESTION]


\end{userprompt}





\subsection{Sub-agents}
\label{app:sub_agents}

% \zixuan{the prompt for IMAS is slightly different due to the nature of the task}

\begin{subagent}[title={CoTAgent}]

\pykw{async} \pykw{def} \pyfn{CoTAgent}(\pykw{self}, agent\_input, model: \pykw{str}):

\quad \pycom{\# Basic setting}

\quad temperature = 0.5

\quad \pycom{\# Chain-of-Thought instruction}

\quad cot\_instruction = \pystr{"Please think step by step and then solve the task."}

\quad \pycom{\# Instantiate CoT LLM}

\quad cot\_agent = LLMAgentBase([\pystr{thinking}, \pystr{answer}],\pystr{Chain-of-Thought LLM}, model=model, temperature=temperature)

\quad thinking, answer = \pykw{await} cot\_agent([agent\_input], cot\_instruction)

\quad final\_answer = self.make\_final\_answer(thinking, answer)

\quad \pykw{return} final\_answer \\

func_string = inspect.getsource(CoTAgent) \\

COT = \code{\{}

\quad \jsonkey{description}: 
\jsonstr{By encouraging the LLM to think step by step rather than directly outputting an answer, chain-of-thought reasoning enables complex problem-solving through intermediate steps.}

\quad \jsonkey{name}: \jsonstr{Chain-of-Thought Agent (CoTAgent)}

\quad \jsonkey{required_arguments}: \code{\{}

\quad\quad \jsonkey{agent_input}: 
\jsonstr{The input for the CoTAgent. If empty, the parser replaces it with the original question.}

\quad \code{\}}

\quad \jsonkey{implementation}: \jsonstr{\{func\_string\}}

\code{\}}


\end{subagent}


\begin{subagent}[title={SCAgent}]
    \pykw{async} \pykw{def} \pyfn{SCAgent}(\pykw{self}, agent\_input, model: \pykw{str}):

\quad \pycom{\# Basic setting}

\quad temperature = 0.5

\quad num\_repeated\_samples = 5

\quad \pycom{\# Chain-of-Thought instruction}

\quad cot\_instruction = \pystr{"Please think step by step and then solve the task."}

\quad \pycom{\# Initialize multiple CoT agents for self-consistency}

\quad cot\_agents = [
\quad\quad LLMAgentBase([\pystr{thinking}, \pystr{answer}],
\quad\quad \pystr{Chain-of-Thought LLM},
\quad\quad model=model,
\quad\quad temperature=temperature)
\quad\quad \pykw{for} \_ \pykw{in} range(num\_repeated\_samples)
\quad ]

\quad thinking\_mapping = \{\}

\quad answer\_mapping = \{\}

\quad possible\_answers = []

\quad \pykw{for} i \pykw{in} range(num\_repeated\_samples):

\quad\quad thinking, answer = \pykw{await} cot\_agents[i]([agent\_input], cot\_instruction)

\quad\quad possible\_answers.append(answer.content)

\quad\quad thinking\_mapping[answer.content] = thinking

\quad\quad answer\_mapping[answer.content] = answer

\quad \pycom{\# Ensembling answers via majority voting}

\quad answer = self.majority\_voting(possible\_answers)

\quad thinking = thinking\_mapping[answer]

\quad answer = answer\_mapping[answer]

\quad final\_answer = self.make\_final\_answer(thinking, answer)

\quad \pykw{return} final\_answer \\

func\_string = inspect.getsource(SCAgent) \\

COT\_SC = \code{\{}

\quad \jsonkey{description}:
\jsonstr{While an LLM can arrive at the correct answer, its reasoning may vary. By repeatedly asking the same question with higher temperature settings, multiple reasoning paths are generated. These answers from Chain-of-Thought agents are then combined through ensembling to produce a more accurate final answer. This approach is best suited for problems requiring high confidence through consensus.}

\quad \jsonkey{name}:
\jsonstr{Self-Consistency with Chain-of-Thought (SCAgent)}

\quad \jsonkey{required_arguments}: \code{\{}

\quad\quad \jsonkey{agent_input}:
\jsonstr{The input for the SCAgent. If empty (""), the parser automatically replaces it with the original question.}

\quad \code{\}}

\quad \jsonkey{implementation}:
\jsonstr{\{func\_string\}}

\code{\}}

\end{subagent}


\begin{subagent}[title={DebateAgent}]
\pykw{async} \pykw{def} \pyfn{DebateAgent}(\pykw{self}, agent\_input, model: \pykw{str}, debate\_roles: List[\pykw{str}]):

\quad \pycom{\# Basic setting}

\quad temperature = 0.5

\quad max\_debate\_round = 5

\quad \pycom{\# Instruction for initial reasoning}

\quad debate\_initial\_instruction = 
\quad \pystr{"Please think step by step and then solve the task."}

\quad \pycom{\# Instruction for debate updates}

\quad debate\_instruction = 
\quad \pystr{"Given solutions from other agents, consider their opinions as advice and provide an updated answer."}

\quad \pycom{\# Initialize debate agents with different roles}

\quad debate\_agents = [
\quad\quad LLMAgentBase([\pystr{thinking}, \pystr{answer}],
\quad\quad \pystr{Debate LLM},
\quad\quad model=model,
\quad\quad role=role,
\quad\quad temperature=temperature)
\quad\quad \pykw{for} role \pykw{in} debate\_roles
\quad ]

\quad \pycom{\# Instruction for final decision}

\quad final\_decision\_instruction =
\quad \pystr{"Given all reasoning and answers, carefully provide a final answer."}

\quad final\_decision\_agent = LLMAgentBase(
\quad\quad [\pystr{thinking}, \pystr{answer}],
\quad\quad \pystr{Final Decision LLM},
\quad\quad model=model,
\quad\quad temperature=temperature
\quad )

\quad all\_thinking = [[] \pykw{for} \_ \pykw{in} range(max\_debate\_round)]

\quad all\_answer = [[] \pykw{for} \_ \pykw{in} range(max\_debate\_round)]

\quad \pycom{\# Perform debate rounds}

\quad \pykw{for} r \pykw{in} range(max\_debate\_round):

\quad\quad \pykw{for} i \pykw{in} range(len(debate\_agents)):

\quad\quad\quad \pykw{if} r == 0:

\quad\quad\quad\quad thinking, answer = 
\quad\quad\quad\quad \pykw{await} debate\_agents[i]([agent\_input], debate\_initial\_instruction)

\quad\quad\quad \pykw{else}:

\quad\quad\quad\quad input\_infos = 
\quad\quad\quad\quad [agent\_input] 
\quad\quad\quad\quad + [all\_thinking[r-1][i]] 
\quad\quad\quad\quad + all\_thinking[r-1][:i] 
\quad\quad\quad\quad + all\_thinking[r-1][i+1:]

\quad\quad\quad\quad thinking, answer = 
\quad\quad\quad\quad \pykw{await} debate\_agents[i](input\_infos, debate\_instruction)

\quad\quad\quad all\_thinking[r].append(thinking)

\quad\quad\quad all\_answer[r].append(answer)

\quad \pycom{\# Final decision based on all debate results}

\quad thinking, answer = 
\quad \pykw{await} final\_decision\_agent(
\quad\quad [agent\_input] 
\quad\quad + all\_thinking[max\_debate\_round-1] 
\quad\quad + all\_answer[max\_debate\_round-1],
\quad\quad final\_decision\_instruction
\quad )

\quad final\_answer = self.make\_final\_answer(thinking, answer)

\quad \pykw{return} final\_answer \\

func\_string = inspect.getsource(DebateAgent) \\

LLM\_debate = \code{\{}

\quad \jsonkey{description}:
\jsonstr{By letting different LLMs debate with each other, diverse perspectives are leveraged to reach better solutions. This agent is best suited for problems that benefit from multiple viewpoints.}

\quad \jsonkey{name}:
\jsonstr{LLM Debate (DebateAgent)}

\quad \jsonkey{required_arguments}: \code{\{}

\quad\quad \jsonkey{agent_input}:
\jsonstr{The input for the DebateAgent. If empty (""), the parser automatically replaces it with the original question.}

\quad\quad \jsonkey{debate_roles}:
\jsonstr{A list of roles (must include more than one) representing distinct perspectives, such as "Mathematics Professor" or "Statistician".}

\quad \code{\}}

\quad \jsonkey{implementation}:
\jsonstr{\{func\_string\}}

\code{\}}

\end{subagent}


\begin{subagent}[title=ReflexionAgent]
\pykw{async} \pykw{def} \pyfn{ReflexionAgent}(\pykw{self}, agent\_input, model: \pykw{str}):

\quad \pycom{\# Basic setting}

\quad temperature = 0.5

\quad max\_reflection\_round = 5

\quad \pycom{\# Instruction for initial reasoning}

\quad initial\_instruction = 
\quad \pystr{"Please think step by step and then solve the task."}

\quad \pycom{\# Instruction for reflection and refinement}

\quad reflect\_instruction = 
\quad \pystr{"Given previous attempts and feedback, carefully consider where you could go wrong in your latest attempt. Using insights from previous attempts, try to solve the task better."}

\quad cot\_agent = LLMAgentBase(
\quad\quad [\pystr{thinking}, \pystr{answer}],
\quad\quad \pystr{Chain-of-Thought LLM},
\quad\quad model=model,
\quad\quad temperature=temperature
\quad )

\quad \pycom{\# Instruction for critic feedback}

\quad critic\_instruction =
\quad \pystr{"Please review the answer above and criticize where it might be wrong. If you are absolutely sure it is correct, output exactly 'True' in 'correct'."}

\quad critic\_agent = LLMAgentBase(
\quad\quad [\pystr{feedback}, \pystr{correct}],
\quad\quad \pystr{Critic LLM},
\quad\quad model=model,
\quad\quad temperature=temperature
\quad )

\quad \pycom{\# Initial attempt}

\quad cot\_inputs = [agent\_input]

\quad thinking, answer = 
\quad \pykw{await} cot\_agent(cot\_inputs, initial\_instruction, 0)

\quad \pycom{\# Iterative self-reflection loop}

\quad \pykw{for} i \pykw{in} range(max\_reflection\_round):

\quad\quad feedback, correct = 
\quad\quad \pykw{await} critic\_agent(
\quad\quad\quad [agent\_input, thinking, answer],
\quad\quad\quad critic\_instruction,
\quad\quad\quad i
\quad\quad )

\quad\quad \pykw{if} correct.content == \pystr{'True'}:

\quad\quad\quad \pykw{break}

\quad\quad cot\_inputs.extend([thinking, answer, feedback])

\quad\quad thinking, answer =
\quad\quad \pykw{await} cot\_agent(cot\_inputs, reflect\_instruction, i + 1)

\quad final\_answer = self.make\_final\_answer(thinking, answer)

\quad \pykw{return} final\_answer \\

func\_string = inspect.getsource(ReflexionAgent) \\

Reflexion = \code{\{}

\quad \jsonkey{description}:
\jsonstr{To enhance performance, an LLM can iteratively improve its answer based on feedback. By reflecting on previous attempts and incorporating critique, the model refines its reasoning and produces a more accurate solution. This agent is best suited for complex problems that benefit from self-correction.}

\quad \jsonkey{name}:
\jsonstr{Self-Refine (Reflexion)}

\quad \jsonkey{required_arguments}: \code{\{}

\quad\quad \jsonkey{agent_input}:
\jsonstr{The input for the ReflexionAgent. If empty (""), the parser automatically replaces it with the original question.}

\quad \code{\}}

\quad \jsonkey{implementation}:
\jsonstr{\{func\_string\}}

\code{\}}


\end{subagent}

    
